\unnumberedChapter{The Cab Handheld}

Now that we have a base station, we can enter commands and control locomotives. However, this is not particularly comfortable and can actually be quite cumbersome. What we need is a device we can hold in our hands, with familiar knobs and buttons. In other words, we need a cab handheld.

Cab handhelds are primarily used to control locomotives. Depending on their additional capabilities, they may also allow the user to configure decoders or operate turnouts. Handhelds are typically connected to the bus, for example LocoNet, although sometimes a separate bus is provided specifically for handhelds. Traditionally, a cable connects the handheld to access points distributed around the layout.

Just as with base stations and boosters, there is no shortage of commercially available cab handhelds. More recently, wireless handhelds have become increasingly popular. And, of course, some base stations integrate handheld controls directly into their front panels.

This chapter describes a general-purpose handheld designed specifically for locomotive control. It connects directly to the LCS bus via cable and provides the basic controls for selecting a locomotive, setting its speed and direction, and operating its function keys.

Implementing a base station and a handheld is all that is needed to run a locomotive and finally see something moving as a result of all the hard work involved in building a layout control system so far.
The cab handheld described here is initially designed as a development board for the firmware. Nevertheless, it can also be used as a fully functional cab handheld. Later versions will build upon the same firmware but use a more convenient form factor.

\unnumberedSection{Requirements}

A cab handheld needs to provide all the basic controls required to operate a locomotive. This also implies having local non-volatile memory to remember previously selected locomotives. This makes it easy to switch between a small collection of locomotives while retaining their individual characteristics.
A display shows the current state of the handheld and, together with the configuration buttons, allows the user to configure its operation.

Looking at commercially available handhelds, they often seem to resemble TV remote controls. They typically feature a numeric keypad, a collection of up and down buttons, and a speed control knob. (No offense intended!) In all fairness, these devices are designed to control not only locomotives but also various other aspects of the layout.
But what about a cab handheld that focuses exclusively on locomotive operation, rather than trying to provide all the controls needed to operate an entire layout?

Our cab handheld takes a different approach. Instead of relying on a large collection of generic buttons, it provides dedicated controls for functions such as a horn or whistle, a bell, and so on. There are also configuration buttons, dedicated direction controls, and a small set of general-purpose buttons that can be assigned to locomotive-specific functions.

And, of course, there is the rotary knob for setting the locomotive speed.
The following figure shows a rough sketch of the proposed cab handheld and its controls.

\begin{center}
	\begin{tikzpicture}[scale=0.9, transform shape]
        
     	% \draw[help lines, gray!50, dashed] (0,0) grid(9,14);
     	
     	\node[ tsRoundedRectangle, 
                minimum width=9cm,
                minimum height=14cm,
                text width=3cm,
                text centered,
                fill=white!50] (display) at (4.5,7) {};

       	\node[ tsRoundedRectangle, 
                minimum width=3.5cm,
                minimum height=3.5cm,
                text width=3cm,
                text centered,
                fill=gray!30] (display) at (4.5,9.5) {Display};
        
        \node[ tsRoundedRectangle, 
                minimum width=1.2cm,
                minimum height=1cm,
                text width=1.2cm,
                text centered,
                fill=red!50] (horn) at (6.5,12.5) {Horn};
        
        \node[ tsRoundedRectangle, 
                minimum width=1cm,
                minimum height=1cm,
                text width=1cm,
                text centered,
                fill=blue!20] (menu) at (1.5,10.5) {Men};
                
      	\node[ tsRoundedRectangle, 
                minimum width=1cm,
                minimum height=1cm,
                text width=1cm,
                text centered,
                fill=blue!20] (up) at (7.5,10.5) {Up};

     	\node[ tsRoundedRectangle, 
                minimum width=1cm,
                minimum height=1cm,
                text width=1cm,
                text centered,
                fill=blue!20] (sel) at (1.5,8.5) {Sel};
                
    	\node[ tsRoundedRectangle, 
                minimum width=1cm,
                minimum height=1cm,
                text width=1cm,
                text centered,
                fill=blue!20] (down) at (7.5,8.5) {Dn};
                
     	\node[ tsRoundedRectangle, 
                minimum width=1cm,
                minimum height=1cm,
                text width=1cm,
                text centered,
                fill=red!50] (f1) at (1.5,6) {F1};
                
      	\node[ tsRoundedRectangle, 
                minimum width=1cm,
                minimum height=1cm,
                text width=1cm,
                text centered,
                fill=red!50] (f2) at (3.5,6) {F2};
                
      	\node[ tsRoundedRectangle, 
                minimum width=1cm,
                minimum height=1cm,
                text width=1cm,
                text centered,
                fill=red!50] (rev) at (1.5,4) {Rev};
                
      	\node[ tsRoundedRectangle, 
                minimum width=1cm,
                minimum height=1cm,
                text width=1cm,
                text centered,
                fill=red!50] (fwd) at (3.5,4) {Fwd};
                
      	\node[ tsRoundedRectangle, 
                minimum width=1cm,
                minimum height=1cm,
                text width=1cm,
                text centered,
                fill=red!50] (f3) at (5.5,6) {F3};
                
     	\node[ tsRoundedRectangle, 
                minimum width=1cm,
                minimum height=1cm,
                text width=1cm,
                text centered,
                fill=red!50] (f4) at (7.5,6) {F4};
                
       	\node[ tsRoundedRectangle, 
                minimum width=1cm,
                minimum height=1cm,
                text width=1cm,
                text centered,
                fill=red!50] (Bell) at (1.5,1) {Bell};
                
      	\node[ tsCircle,
         		minimum width=3cm,
                minimum height=3cm,
                text width=2cm,
                text centered,
                fill=red!50] (speed) at (6.5,2.5) {Speed};
         
    \end{tikzpicture}
\end{center}

Configuration and part of the normal operation are handled through four buttons surrounding the display. The \texttt{MENU} button allows the user to cycle through the available menus, while the \texttt{SELECT} button selects the currently displayed menu.
This navigation scheme can be nested. Once inside a menu, the \texttt{MENU}, \texttt{UP}, and \texttt{DOWN} buttons are used according to the particular screen, while \texttt{SELECT} typically confirms the selected action.

The direction buttons \texttt{REV} and \texttt{FWD}, together with the \texttt{SPEED} knob, control the speed and direction of the locomotive or consist. \texttt{F1} through \texttt{F4} are four general-purpose buttons that can be mapped to special functions of the particular locomotive. The dedicated \texttt{HORN} and \texttt{BELL} buttons complete the initial design.

The display itself follows a common structure for all screens. The corner fields correspond to the navigation buttons located next to them. The text shown in these fields depends on the current screen and typically indicates the action associated with each button.
Between the two corner fields at the top and bottom of the display, there is a data field containing general status information. Finally, two context-dependent lines in the center of the display provide the main information for the current screen.

As always, there are many ways to build a cab handheld. Although this version connects to the LCS bus via cable, implementing a wireless version would not be particularly difficult. Having more or fewer buttons, or providing a numeric keypad-style input, are all perfectly valid alternatives. Ultimately, it comes down to personal preference.

The current handheld does not provide any controls for accessories such as turnouts. This is a subject better left to dedicated layout control panels and control software. Our cab handheld deliberately favors the idea of a locomotive control stand rather than a TV remote-style device. Since this is ultimately a matter of taste and preference, feel free to build your own.

There is, of course, also the option of connecting commercially available handhelds. This would require a gateway translating, for example, between a LocoNet-based handheld protocol and the LCS protocol.
Well, perhaps one day you will be able to connect such handhelds via the LocoNet bus. For now, however, this remains on the backlog. It is a subject for the Gateway part of this book.